\subsection{对初始条件的依赖性}

设 \(\mathbf{x}^{\prime}=\mathbf{f}(t,\mathbf{x})\) 是 \(I\times H\) 上的局部利普希茨方程；由定理 2（IV，第 172 页），对 \((t_{0},\mathbf{x}_{0})\) 的每个点 \(I\times H\) ，存在最大的区间 \(J(t_{0},\mathbf{x}_{0})\subset I\) （不退化为一点、包含 \(t_{0}\) ），在该区间上存在此方程的一个（且仅有的一个）积分，它等于 \(x_{0}\) 于 \(t_{0}\) ；我们将阐明这个积分以及其定义区间 \(J(t_{0},\mathbf{x}_{0})\) 依赖点 \((t_{0},\mathbf{x}_{0})\) 的方式。

定理 4. 设 \(\mathbf{f}\) 是 \(\mathrm{I} \times \mathrm{H}\) 上的局部利普希茨函数，而 \((a, \mathbf{b})\) 是 \(\mathrm{I} \times \mathrm{H}\) 的任意一点。

\(1^{\circ}\) 存在一个区间 \(\mathbf{K} \subset \mathbf{I}\) （即 \(a\) 在 \(\mathbf{I}\) 中的邻域）和一个 \(\mathbf{V}\) 在 \(\mathbf{b}\) 中的邻域 \(\mathbf{H}\) ，使得对 \((t_0, \mathbf{x}_0)\) 的每个点 \(\mathbf{K} \times \mathbf{V}\) ，存在一个（且仅有一个）积分 \(\mathbf{u}(t, t_0, \mathbf{x}_0)\) ，它定义在 \(\mathbf{K}\) 上、取值于 \(\mathbf{H}\) 中且等于 \(\mathbf{x}_0\) 于点 \(t_0\) （换言之，\(J(t_0, \mathbf{x}_0) \supset \mathbf{K}\) 对一切 \((t_0, \mathbf{x}_0) \in \mathbf{K} \times \mathbf{V}\) 成立）。

\(2^{\circ}\) 映射 \((t,t_0,\mathbf{x}_0)\mapsto \mathbf{u}(t,t_0,\mathbf{x}_0)\)（从 \(\mathrm{K}\times \mathrm{K}\times \mathrm{V}\) 到 H 中）是一致连续的。

\(3^{\circ}\) 存在 \(\mathbf{W} \subset \mathbf{V}\) 在 \(\mathbf{b}\) 中的一个邻域 \(\mathbf{H}\) ，使得对每个点

\[
(t, t _ {0}, \mathbf {x} _ {0}) \in \mathrm{K} \times \mathrm{K} \times \mathrm{W},
\]

方程 \(\mathbf{x}_{0} = \mathbf{u}(t_{0}, t, \mathbf{x})\) 在 V 上有唯一的解 x，它等于 \(\mathbf{u}(t, t_{0}, \mathbf{x}_{0})\)（“把积分关于积分常数解出”）。

\(1^{\circ}\) 设 S 是以 \(\mathbf{b}\) 为中心、以 \(r\) 为半径且包含于 H 中的球，并设 \(\mathrm{J}_0\) 是包含于 I 的区间，即 \(a\) 在 I 中的邻域，使得 \(\mathbf{f}\) 有界且是利普希茨的（关于某个常数 \(k\) ）于 \(\mathrm{J}_0 \times \mathrm{S}\) 上；用 M 表示 \(\| \mathbf{f}(t, \mathbf{x})\|\) 在 \(\mathrm{J}_0 \times \mathrm{S}\) 上的上确界。于是存在（IV，第 171 页，定理 1）一个区间 \(\mathrm{J} \subset \mathrm{J}_0\) （即 \(a\) 在 I 中的邻域）和一个积分 \(\mathbf{v}\) ，即方程 \(\mathbf{x}' = \mathbf{f}(t, \mathbf{x})\) 的一个积分，它定义在 J 上、取值于 S 中且等于 \(\mathbf{b}\) 于 \(a\) 。我们将看到，以 \(\mathbf{b}\) 为中心、以 \(r/2\) 为半径的开球 V，以及 J 与区间 \([a - l, a + l[\) （其中 \(l\) 足够小）的交 K，都合乎要求。的确，IV 第 174 页的命题 7（应用于集合 \(\mathrm{J}_0 \times \mathrm{S}\) 和 \(\alpha = 0\) 的情形）表明存在 \(\mathbf{x}' = \mathbf{f}(t, \mathbf{x})\) 的一个积分，它定义在 K 上、取值于 S 中且等于 \(\mathbf{x}_0\) 于一点 \(t_0 \in \mathrm{K}\) ，其条件是

\[
\left\| \mathbf {v} (t) - \mathbf {b} \right\| + \left\| \mathbf {v} (t _ {0}) - \mathbf {x} _ {0} \right\| e ^ {k | t - t _ {0} |} <   r\tag{18}
\]

对每个 \(t \in \mathbf{K}\) 。于是由中值定理有

\[
\| \mathbf {v} (t) - \mathbf {b} \| \leqslant \mathrm{M} | t - a | \leqslant \mathrm{M} l
\]

对每个 \(t \in \mathbf{K}\) ；由于 \(\| \mathbf{x}_0 - \mathbf{b} \| < r / 2\) ，可见只需取 \(l\) 使得

\[
\mathbf {M} l + (\mathbf {M} l + r / 2) e ^ {2 k l} <   r\tag{19}
\]

便可使关系式 (18) 对每个 \((t, t_{0}, \mathbf{x}_{0})\)（属于 \(K \times K \times V\) ）都成立。

\(2^{\circ}\) 由中值定理有

\[
\left\| \mathbf {u} (t _ {1}, t _ {0}, \mathbf {x} _ {0}) - \mathbf {u} (t _ {2}, t _ {0}, \mathbf {x} _ {0}) \right\| \leqslant \mathrm{M} \left| t _ {2} - t _ {1} \right|\tag{20}
\]

对 K 中的一切 \(t_0, t_1, t_2\) 与 V 中的 \(\mathbf{x}_0\) 。现在命题 5（IV，第 169 页）表明

\[
\left\| \mathbf {u} (t, t _ {0}, \mathbf {x} _ {1}) - \mathbf {u} (t, t _ {0}, \mathbf {x} _ {2}) \right\| \leqslant e ^ {2 k t} \left\| \mathbf {x} _ {2} - \mathbf {x} _ {1} \right\|\tag{21}
\]

对 K 中的每个 t 和 \(t_{0}\)，以及 V 中的 \(x_{1}\) 与 \(x_{2}\)。最后，若 \(t_{1}\) 与 \(t_{2}\) 是 K 中任意两点，

\[
\left\| \mathbf {u} \left(t _ {1}, t _ {2}, \mathbf {x} _ {0}\right) - \mathbf {u} \left(t _ {2}, t _ {2}, \mathbf {x} _ {0}\right) \right\| \leqslant \mathrm{M} \left| t _ {2} - t _ {1} \right|
\]

这是由中值定理得到的，也就是说

\[
\left\| \mathbf {u} (t _ {1}, t _ {2}, \mathbf {x} _ {0}) - \mathbf {x} _ {0} \right\| \leqslant \mathbf {M} \left| t _ {2} - t _ {1} \right|;
\]

由于 \(\mathbf{u}(t, t_2, \mathbf{x}_0)\) 与取值 \(\mathbf{u}(t_1, t_2, \mathbf{x}_0)\) 于点 \(t_1\) 的积分相同，命题 5（IV，第 169 页）表明

\[
\left\| \mathbf {u} (t, t _ {1}, \mathbf {x} _ {0}) - \mathbf {u} (t, t _ {2}, \mathbf {x} _ {0}) \right\| \leqslant \mathrm{Me} ^ {2 k l} | t _ {2} - t _ {1} |\tag{22}
\]

对 \(t, t_{1}, t_{2}\) 中的一切 \(\mathbf{K}\) 与 \(\mathbf{x}_{0}\) 中的一切 \(\mathrm{V}\) 都成立。于是三个不等式 (20)、(21) 与 (22) 证明了映射 \((t, t_{0}, \mathbf{x}_{0}) \mapsto \mathbf{u}(t, t_{0}, \mathbf{x}_{0})\) 在 \(\mathrm{K} \times \mathrm{K} \times \mathrm{V}\) 上的一致连续性。

\(3^{\circ}\) 由 (20)，有 \(\|\mathbf{u}(t, t_{0}, \mathbf{x}_{0}) - \mathbf{x}_{0}\| \leqslant \mathbf{M} |t - t_{0}| \leqslant 2\mathbf{M}l\) 于

\[
\mathbf {K} \times \mathbf {K} \times \mathbf {V}.
\]

若把 \(l\) 取得足够小，使 \(2\mathrm{M}l < r / 4\) ，则可以看出：若 \(\mathbf{x}_0\) 是以 \(\mathbf{W}\) 为中心、以 \(\mathbf{b}\) 为半径的开球 \(r / 4\) 中的任意一点，那么 \(\mathbf{u}(t,t_0,\mathbf{x}_0)\in \mathbf{V}\) 对任何 \(t\) 与 \(t_0\)（都属于 \(\mathbf{K}\) ）成立。若 \(\mathbf{x} = \mathbf{u}(t,t_0,\mathbf{x}_0)\) ，则函数 \(s\mapsto \mathbf{u}(s,t,\mathbf{x})\) 便定义在 \(\mathbf{K}\) 上，并且等于 (1) 的取值 \(\mathbf{x}\) 于点 \(t\) 的那个积分，也就是等于 \(\mathbf{u}(s,t_0,\mathbf{x}_0)\) ；特别地

\[
\mathbf {x} _ {0} = \mathbf {u} (t _ {0}, t _ {0}, \mathbf {x} _ {0}) = \mathbf {u} (t _ {0}, t, \mathbf {x}).
\]

此外，若 \(y \in V\) 满足 \(\mathbf{x}_{0} = \mathbf{u}(t_{0}, t, \mathbf{y})\) ，则积分 \(s \mapsto \mathbf{u}(s, t, \mathbf{y})\) 取值 \(x_{0}\) 于 \(t_{0}\) ，故与 \(s \mapsto \mathbf{u}(s, t_{0}, \mathbf{x}_{0})\) 相同，从而后者在 t 处取值 x，这就表明 y = x 并完成了证明。

